\chapter{Signal storage and the dependency graph}
\label{ch:signalgraph}

\begin{aside}
A signal is a value with a name and a list of subscribers. The
list of subscribers is built automatically by tracking which
effects read which signals. The graph is the runtime's nervous
system.
\end{aside}

\section{The Signal type}

A signal stores a value of any type plus metadata:

\begin{lstlisting}
template<typename T>
class Signal {
    T value_;
    std::vector<EffectId> subs_;
    uint32_t version_;
public:
    const T& get();   // tracks dependency
    void set(T v);    // marks subs dirty
};
\end{lstlisting}

\section{Tracking dependencies}

When an effect runs, it sets a thread-local
``currently-running effect''. Signal reads, while that
thread-local is non-null, add the effect to the signal's
subscriber list and remember the signal in the effect's
dependency list.

After the effect finishes, the thread-local is cleared. The
graph now has the edges this run produced.

\section{Re-tracking each run}

Every time an effect re-runs, its dependencies are tracked
fresh. The previous edges are dropped first, then the new
run accumulates new ones.

This handles conditional reads: an effect that reads
\code{signalA} when X is true and \code{signalB} when X is
false subscribes only to whichever it read this run.

\section{Versioning}

Each signal carries a version counter. \code{set} increments
it. An effect that recomputes can compare the versions of
its dependencies to detect ``did anything actually change''.

This is useful for batching: if many signals update but
their net effect on a derived value is zero, the derived
effect can early-out.

\section{Memos}

A memo is a computed signal: a function of other signals,
cached. It re-runs only when its inputs change.

\begin{lstlisting}
auto fullname = memo([&] {
    return ctx.get(first) + " " + ctx.get(last);
});
\end{lstlisting}

The memo subscribes to \code{first} and \code{last}. Its
own value is itself a signal that other effects can read.

\section{Effects vs memos}

Both react to signal changes. The difference:

\begin{itemize}[leftmargin=*]
  \item \textbf{Memo:} produces a value. Lazy --- runs when
    its value is read.
  \item \textbf{Effect:} produces side effects. Eager --- runs
    when any input changes.
\end{itemize}

A memo that nobody reads does no work. An effect runs
regardless.

\section{Batching updates}

A flurry of signal updates within one event handler should
fire effects once, not N times. \maya{} batches:

\begin{lstlisting}
ctx.batch([&] {
    ctx.set(a, 1);
    ctx.set(b, 2);
    ctx.set(c, 3);
});  // one re-render at the end
\end{lstlisting}

Inside the batch, dirty signals accumulate; effects are
fired after the batch closes.

\section{Cycle detection}

An effect that writes a signal it reads can loop forever.
\maya{} detects this in debug builds: if an effect's own
write triggers its own re-run, abort with a diagnostic.

In release, the cycle still happens; the user sees a
runaway CPU.

\section{Graph size}

A small app has a few dozen signals and a few dozen
effects. The graph fits in 10 KB. Lookups are
\code{vector<EffectId>}; insertions are amortised O(1).

\begin{warning}
A signal with thousands of subscribers (everyone in the UI
reads a global ``current time'') is fine but expensive. Each
update walks every subscriber. Prefer narrow signals.
\end{warning}

\section{Recap}

\begin{recap}
\begin{itemize}[leftmargin=*]
  \item Signals: value + version + subscribers.
  \item Effects track dependencies via a thread-local.
  \item Memos are lazy; effects are eager.
  \item Batch to coalesce updates.
  \item Cycles caught in debug; expensive in release.
\end{itemize}
\end{recap}

\section*{Workshop}

\begin{enumerate}[leftmargin=*]
  \item Build a counter signal. Subscribe an effect that
    logs each change.
  \item Create a memo that depends on two signals.
    Confirm it re-runs only when either changes.
  \item Try to create a cycle (effect writes its own
    read). Observe the diagnostic.
\end{enumerate}
